\section{Budgets and controlled measurement}
\label{sec:budgets}

\subsection{Admission and accounting}
\label{sec:budgets-admission}

A runtime needs two mechanisms to limit spending. Before a model or tool call, it must decide
whether the call may start. After execution, it must record the cost. The first mechanism uses
information available before the call; the second uses the usage reported afterwards.

\begin{definition}[Admission]
\label{def:admission}
Admission decides whether a unit of work may start. It checks the run's existing commitments
against a limit, then permits or refuses the work. Refused work never runs and incurs no charge.
\end{definition}

\begin{definition}[Accounting]
\label{def:accounting}
Accounting records the cost of a unit of work after execution. It adds the incurred charge to a
running total, using the same units as the limit. Admission reads this total.
\end{definition}

Admission determines whether work may start; accounting records its cost. Each depends on the
other. If accounting omits a call, the admission check can permit work that the limit would
otherwise exclude. If no admission check reads the total, accounting merely records an overrun
after it occurs. \textbf{Accounting maintains the usage total; admission uses it to enforce the
limit.}

Accounting reads the charge reported with each reply. Admission must decide before that reply
exists. It knows the input to be sent but does not yet know the output. This difference makes a
call-count limit easier to enforce than a token limit: each call counts as one, while its token
cost becomes known only after execution.

Our planning agent uses call-count admission. It checks a budget of forty calls before each model
call and records the call's input and output tokens. The token admission and reservation mechanism
developed below is a proposed design. Our code does not implement it.

Call counts measure usage coarsely. In our two recorded runs, a single call's input ranges from
571 to 13{,}550 tokens because each controller call resends the growing conversation. A budget of
forty calls therefore bounds the number of replies. The resulting bound on token cost depends
on the largest possible call.

\subsection{A budget invariant}
\label{sec:budgets-invariant}

Admission and accounting need an explicit budget unit. The unit determines which cost the budget
limits. We consider three choices.

\bi

\item \textbf{Output tokens.} Count only tokens generated by the model. A per-call output limit
bounds each call's contribution in advance, simplifying admission. This budget excludes input
tokens, whose count grows as calls resend the conversation.

\item \textbf{All processed tokens.} Count input and output tokens together to track the provider's
work. Admission requires bounds on both the input and output of each call.

\item \textbf{Money.} Weight tokens by their prices. Token kinds can have different prices, as with
cache reads and fresh input. Accounting must therefore retain separate counts until it computes
the charge. Price changes alter how much work the budget permits.

\ei

\noindent The following argument applies to all three units. The unit determines how we calculate
the charge $\gamma$ of a call.

We express the budget as an invariant over three quantities. $B$ is the budget. $S$ is the settled
spend, the sum of recorded charges for completed calls. $R$ is the sum of outstanding reservations,
the amounts set aside for admitted calls that have not yet been settled. The invariant is
\[
  S + R \;\leq\; B .
\]
A \emph{reservation} $\rho$ is the amount set aside for a call before it starts. It must be a
justified upper bound on that call's cost. \emph{Settlement} replaces the reservation with the
recorded charge $\gamma$ after the call completes.

The required property concerns the charges themselves: the total charge for all calls started by
the run, including unfinished calls, must never exceed $B$. As with the call counter in
\cref{sec:state-invariants}, the runtime can compare only recorded values. To connect the record's
invariant to the charges, we require
\[
  \forall i.\; \gamma_i \leq \rho_i ,
\]
where $i$ ranges over admitted calls, $\rho_i$ is the reservation for call $i$, and $\gamma_i$ is
its charge on completion. This obligation constrains how admission computes $\rho$. It is not a
state predicate: the final charge $\gamma_i$ of an outstanding call is not yet known to the
runtime. The input and output bounds below justify this obligation. If the obligation holds,
$S + R$ bounds the total charge of all calls started, and the invariant implies the required
property. We use the invariant method (\cref{def:invariant}) to prove that the inequality holds
initially and that each transition preserves it.

\bi

\item \textbf{Initially.} No call has been admitted or settled, so $S = 0$ and $R = 0$. The invariant
therefore holds for any $B \geq 0$.

\item \textbf{Admit.} If $S + R + \rho \leq B$, admission reserves $\rho$ for the call by setting
$R$ to $R + \rho$. The guard ensures that the invariant holds after this step. Otherwise,
admission refuses the call and leaves the record unchanged.

\item \textbf{Settle.} The call completes with charge $\gamma$. Settlement sets $S$ to $S + \gamma$
and $R$ to $R - \rho$, changing their sum by $\gamma - \rho$. Because $\gamma \leq \rho$, the sum
cannot grow, and $S + R \leq B$ still holds. This step also releases the unused amount
$\rho - \gamma$ to the unreserved budget.

\item \textbf{Release.} If an admitted call is never issued, for example because the run ends
first, the runtime removes its reservation and reduces $R$ by $\rho$. The call incurred no charge,
so the invariant holds.

\item \textbf{Every other transition.} Leave $S$ and $R$ unchanged. In particular, no transition
decreases $S$. Settled spend has the run lifetime defined in \cref{sec:state-lifetimes}. Resetting
it to zero would make the budget available again while leaving the charges incurred.

\ei

Concurrent calls require one further condition: the guard check and reservation must form a
single step. Suppose two model calls are admitted in the same superstep
(\cref{sec:boundary-scheduling}). If both check $S + R + \rho \leq B$ before either adds its
reservation to $R$, both can pass while their combined cost exceeds $B$.

The proof depends on establishing $\gamma \leq \rho$. For a budget measured in all processed
tokens, $\rho$ is the sum of an input bound and an output bound. Each needs justification.

\bi

\item \textbf{Input.} Establish an upper bound on the call's input tokens. The provider's counting
endpoint returns an estimate before the call is sent: ``The token
count is an estimate. In some cases, the actual number of input tokens used when creating a message
might differ by a small amount''~\citep{anthropic-token-counting}. Comparing past estimates with
recorded charges can inform a margin. A strict guarantee, however, requires a justified upper
bound on the estimation error. Alternatively, admission can use a limit the request cannot exceed,
such as the model's context window.

\item \textbf{Output.} Use an enforced upper bound on output tokens. The provider enforces the
per-call output limit sent with the request and stops generation when that limit is
reached~\citep{anthropic-stop-reasons}. This bound holds regardless of the content the model
intended to generate.

\ei

An estimate alone cannot justify the invariant. If admission computes $\rho$ from an estimate and
the call costs more, settlement increases the sum by $\gamma - \rho > 0$. The record contains no
flag identifying the step at which $S + R$ exceeded $B$. \textbf{This reservation-based guarantee
requires an upper bound on each call's charge.} Estimates support a forecast of spending. The
claim that spending cannot exceed $B$ requires the bounds above.

The runtime must also handle unused reservations and uncertain charges. Releasing unused
reservations is a liveness requirement: an unreleased reservation keeps $R$ high and causes
admission to refuse work that would otherwise fit. Leaving an unused reservation outstanding
does not violate the spending invariant. An uncertain charge requires a different response.
If a timeout or crash prevents receipt of a reply, the provider may
still have charged for the call. Releasing its reservation could omit that charge from the record.
The reservation must remain outstanding until the runtime reconciles the charge with the
provider's usage records. Only then can it settle the call or release the reservation.

\subsection{Exhaustion and truncation}
\label{sec:budgets-limits}

A budget and a per-call limit stop work at different points. The budget prevents new work from
starting. The per-call limit stops generation within an admitted call. A run can reach either
limit, with a different result in each case.

\bi

\item \textbf{Exhaustion.} When admission refuses the next call, that call never starts. The run
ends with a declared status, and a terminal event in the trace records the reason
(\cref{tab:control-outcomes}). Exhaustion leaves calls already in flight running until they
complete and are settled.

\item \textbf{Truncation.} When a call reaches its output limit, the provider stops generation and
reports the reason in the response. It returns a successful response containing the content
generated before the limit.

\ei

Truncation matters most for tool requests. The provider's documentation describes the problem:
``If Claude's response is cut off because it hit the \texttt{max\_tokens} limit,
and the truncated response contains an incomplete tool use block, you'll need to retry the request
with a higher \texttt{max\_tokens} value to get the full tool use''~\citep{anthropic-stop-reasons}.
A truncated request can still resolve to a known tool and bind to its signature. For example, an
edit request can have valid argument names but incomplete replacement text. The dispatch checks in
\cref{sec:boundary-dispatch} test names, arity and policy; they do not check completeness.
\textbf{The stop reason can flag a potentially incomplete tool request even when it passes every
dispatch check.} The runtime must therefore read the stop reason before dispatch. It must refuse an
incomplete request before any effect occurs and return an observation on which the model can act.

Recovery from truncation charges both the original call and the retry. The truncated call ran and
incurred a charge, so accounting settles it normally. A retry with a larger output limit is a new
call. It requires a separate admission check and a reservation sized for the larger limit.

Admission also applies to calls after the main work. Retries, plan revisions and final assessments
generated by a model each require an admitted model call. If the controller uses the entire
budget, the run has no budget left for assessment and ends exhausted. In our planning agent, the
validator uses the same admission check as the other two model nodes. If a patch passes its checks
after the fortieth call, the run ends \demoref{budget\_exhausted} without a verdict. A design that
requires a final assessment reserves part of the budget for it in advance. Other roles are then
admitted against $B$ minus that allocation. Our code does not implement this reservation.

Accounting must cover every role and attempt. Counting only controller calls would leave the
planner's and validator's spending unlimited. Keeping a budget per attempt would reset it after
each failed attempt. Settled spend must have the run lifetime (\cref{sec:state-lifetimes}), as the
call counter does across the proposed attempt edge in \cref{sec:control-backedges}. In our
planning agent, all three model nodes update a shared call counter and usage total. Each model-call
event in the trace identifies its role, so the trace can reconstruct call counts by role.

\subsection{Controlled measurement}
\label{sec:budgets-measurement}

Comparing the cost of two realizations requires a common budget and a controlled design. The
scripted comparison in \cref{sec:events-frameworks} tests whether two realizations implement the
same transition rules for the scripted replies. It does not measure behaviour with a live model, whose replies depend
on the messages each realization constructs. A live comparison measures the model and harness
together: ``When we evaluate `an agent,' we're evaluating the harness and the model working
together''~\citep{anthropic-evals-agents}. Its unit is the trial.

\begin{definition}[Trial]
\label{def:trial}
``Each attempt at a task is a trial. Because model outputs vary between runs, we run multiple trials
to produce more consistent results.''
\end{definition}

A comparison varies the realization, holds other conditions fixed, and repeats each configuration
over several trials. It holds seven conditions fixed.

\bi

\item \textbf{Tasks.} The same alerts, in the same fixture, for every realization.

\item \textbf{Initial environment.} Each trial starts from a fresh copy of the same workspace.
This prevents access to files, caches or history from earlier trials only if tools cannot read
outside the workspace. Our tools permit such access through the glob argument of
\demoref{search}, as \cref{sec:boundary-demo} shows.

\item \textbf{Model configuration.} The same model, version, sampling parameters and per-call output
limit.

\item \textbf{Prompts.} The same system prompts and task messages for each role.

\item \textbf{Tools.} The same tool set, with the same signatures, descriptions and limits.

\item \textbf{Acceptance rules.} The same checks and the same validator, so that \emph{accepted}
means the same thing in both arms.

\item \textbf{Budget policy.} The same unit, $B$, and admission rule. A realization that makes more
calls reaches the limit sooner. Different budget policies could make a policy effect appear to be
an effect of the realization.

\ei

For each realization, we report the distribution of terminal outcomes across trials, together
with i) usage by token kind, ii) calls by role, iii) latency, and iv) the exhaustion rate. Acceptance
rates alone cannot distinguish rejection from budget exhaustion. Usage alone cannot distinguish
low-cost completion from an early stop.

A fixed seed does not establish equivalent trajectories. It fixes how the sampler draws from a
distribution for a given input. It does not fix that input. If two realizations format messages
differently, the model receives different inputs from the first call. A difference in one tool
result also changes all later messages. Whether identical requests produce identical outputs
depends in part on how the provider serves them, which the seed does not control. Our planning
agent fixes seed 0 and temperature 0. These settings fix the sampler's configuration without
determining the trajectory. Repeated trials make it possible to measure differences between
realizations.

Such a measurement supports conclusions about the tested realizations, model, prompts, tasks and
budget. It does not establish which realization performs better with another model, which may
respond differently to the same messages, or under another budget, which changes where runs stop.
A small number of trials may also be insufficient to separate a real difference from variation
between runs.

\subsection{Reading}

revisit \citet{mit6042-lec16}, \S2 (p.~5), for the proof method used in the reservation
invariant; and read \citet{anthropic-evals-agents}, the sections The structure of an
evaluation, How to think about non-determinism in evaluations for agents, and Design the eval
harness and graders, on trials, the harness, and isolated environments.

For reference while writing HW1: \citet{anthropic-token-counting}, on what the counting endpoint
counts and its warning that the result is an estimate; and \citet{anthropic-stop-reasons}, the
sections on \texttt{max\_tokens}, incomplete tool use blocks, and handling truncated responses.

\paragraph{Looking ahead.} In this account of an agent, a named component makes each decision,
and a named artifact records it. A routing function selects the next location, and a routing event
records the choice. The dispatcher permits or refuses a tool request, and an observation records
the result. Admission permits a call, and accounting records its cost. The runtime sets the
terminal status, and the trace records it. HW1 asks you to describe your agent in the same terms.
For each decision, identify the responsible component and the artifact that provides evidence of
the decision. If either is missing, a reviewer cannot check what occurred.
